\documentclass[12pt]{article}
\usepackage[utf8]{inputenc}
\usepackage[T1]{fontenc}
\usepackage{hyperref}
\title{Pictish Language Resources for K-12 Curriculum Development}
\author{Compiled by Claude Sonnet 5.5 (claude-sonnet-5-5) for Andr\'as Kornai}
\date{}
\begin{document}
\maketitle

\section{Introduction}
Pictish is an extinct language of early medieval Scotland. It has no native speakers and no surviving continuous texts. It is known mainly from a small body of ogham inscriptions, place names and personal names. Its classification is disputed. Some scholars see it as a Brythonic (Celtic) dialect, and the inscriptions remain largely undeciphered. A curriculum that uses ``the native language'' is therefore not possible in the ordinary sense. What exists is scholarly material about the evidence. My search found no school programmes and no teaching materials.

\section{Primary Educational Resources}
I found no primers, textbooks or school programmes for Pictish. The following are scholarly or popular overviews that a teacher could adapt for background reading.
\begin{itemize}
\item Carla Nayland, ``Pictish language: Inscriptions'', an essay surveying the Pictish ogham inscriptions: \url{https://carlanayland.org/essays/pictish_inscriptions.htm}
\item David M. M., ``Linguistic Background to Pictish and the Pictish Symbols Writing System'' (2026 note): \url{https://davidmm.web.unc.edu/2026/03/16/note-1-linguistic-background-to-pictish-and-the-pictish-symbols-writing-system/}
\end{itemize}

\section{Linguistic and Reference Materials}
\begin{itemize}
\item Oliver Padel, \emph{Inscriptions of Pictland} (1972 Edinburgh thesis), with a catalogue of the ogham inscriptions: \url{https://era.ed.ac.uk/handle/1842/9863}
\item The same thesis, full item record: \url{https://era.ed.ac.uk/items/5905762f-5dcc-4a9b-a685-4ee63bb7f34b/full}
\item The same thesis, PDF download: \url{https://era.ed.ac.uk/bitstreams/3e9154fd-a18b-4a62-a399-cdd209387ae7/download}
\item Verbix documentation page on Pictish: \url{https://docs.verbix.com/Languages/Pictish}
\item Academia.edu paper, ``Pictish Language The Mysterious Ogham Inscriptions of Scotland'' (an interpretive, non-consensus attempt): \url{https://www.academia.edu/125488632}
\item Conference presentation abstract noting that the inscriptions remain undeciphered: \url{https://istina.ipmnet.ru/conferences/presentations/404295798/}
\end{itemize}

\section{Religious and Community Resources}
I found no religious texts or community organisations for the Pictish language in my searches.

\section{Digital and Audio Resources}
I found no audio recordings, software or dedicated digital corpus. The online items above are the only digital material I saw.

\section{Curriculum Development Context}
Because the language is unattested as connected text, any K-12 use would be heritage and history oriented. Possible topics are the ogham script, Pictish symbol stones and place-name study. Teachers should present the uncertainty honestly, including the disagreement over whether Pictish was Brythonic, and should not invent vocabulary. Further searching would be needed for museum and heritage-education programmes in Scotland.

\end{document}